


\def\stat{egorova-zac}





\def\tit{МАШИННЫЙ ПЕРЕВОД:\\ ИНДИКАТОРНАЯ ОЦЕНКА РЕЗУЛЬТАТОВ 


ОБУЧЕНИЯ\\ ИСКУССТВЕННОЙ НЕЙРОННОЙ СЕТИ}





\def\titkol{Машинный перевод: индикаторная оценка результатов 


обучения ИНС} %искусственной нейронной сети}








\def\aut{А.\,Ю.~Егорова$^1$, И.\,М.~Зацман$^2$, М.\,Г.~Кружков$^3$, 


В.\,А.~Нуриев$^4$}





\def\autkol{А.\,Ю.~Егорова, И.\,М.~Зацман, М.\,Г.~Кружков, 


В.\,А.~Нуриев}





\titel{\tit}{\aut}{\autkol}{\titkol}





\index{Егорова А.\,Ю.}


\index{Зацман И.\,М.}


\index{Кружков М.\,Г.}


\index{Нуриев В.\,А.}


\index{Egorova A.\,Yu.}


\index{Zatsman I.\,M.}


\index{Kruzhkov M.\,G.}


\index{Nuriev V.\,A.}








%{\renewcommand{\thefootnote}{\fnsymbol{footnote}}


%\footnotetext[1]


%{Работа выполнена при частичной поддержке РФФИ (проект 18-07-00252) и~в~соответствии с~программой % 


%Московского центра фундаментальной и~прикладной математики.}} 








\renewcommand{\thefootnote}{\arabic{footnote}}


\footnotetext[1]{Институт проблем информатики Федерального 


исследовательского центра <<Информатика и~управление>> 


Российской академии наук, \mbox{ann.shurova@gmail.com}}


\footnotetext[2]{Институт проблем информатики Федерального 


исследовательского центра <<Информатика и~управление>> 


Российской академии наук, \mbox{izatsman@yandex.ru}}


\footnotetext[3]{Институт проблем информатики Федерального 


исследовательского центра <<Информатика и~управление>> 


Российской академии наук, \mbox{magnit75@yandex.ru}}


\footnotetext[4]{Институт проблем информатики Федерального 


исследовательского центра <<Информатика и~управление>> 


Российской академии наук, \mbox{nurieff.v@gmail.com}}








 \label{st\stat}





 \thispagestyle{headings}


 


 \vspace*{-12pt}  








\Abst{Представлены данные, полученные в~ходе наблюдения за обуче\-ни\-ем 


системы нейронного машинного перевода (НМП). Проведена количественная 


оценка работы системы НМП с~помощью индикаторов. В~качестве 


экспериментального материала были использованы 250~русскоязычных 


текстовых фрагментов, для каждого из которых ежемесячно в~течение одного 


года фиксировался перевод на французский язык, выполненный с~помощью 


системы НМП Google. Фиксация переводов была реализована посредством их 


аннотирования в~надкорпусной базе данных (НБД), в~результате чего была 


сформирована серия из 12~аннотаций для каждого из 250~текстовых 


фрагментов. Аннотирование переводов позволило не только зафиксировать 


допущенные в~переводе ошибки в~случае их наличия, но и~определить 


категорию нестабильности НМП, указывающую на изменения качества 


перевода или на их отсутствие. Цель статьи~--- представить разработанный 


индикаторный подход и~пример его применения для оценки результатов 


обучения искусственной нейронной сети (ИНС).}








\KW{нейронный машинный перевод; нестабильность машинного перевода; 


индикаторная оценка; лингвистическое аннотирование; виды нестабильности}





\DOI{10.14357\!/\!08696527200412} 








 \thispagestyle{headings}


 


 \vspace*{-12pt}








\section{Введение}


   


   Компьютерные системы после обучения ИНС


    следуют правилам решения задачи, которые определяются данными 


в~процессе ее обучения (data-driven), а~не создаются заранее человеком. 


Алгоритмы, определяемые данными, не соответствуют правилам, которым 


следуют люди при решении этой же задачи. Согласно Фенстермахеру, 


компьютерные системы могут выполнять ту же задачу, используя 


альтернативные правила и~другие формы представления знания, при-\linebreak


\vspace*{-12pt}





\pagebreak





\noindent


меняемого 


при ее решении~[1]. Например, машинный алгоритм распознавания лиц, 


сформированный в~процессе обучения компьютерной системы, не совпадает 


с~алгоритмом, разработанным человеком~[2].


   


   Основная причина отличий машинных и~разработанных человеком 


алгоритмов заключается в~следующем. Искусственная нейронная сеть


 состоит из слоев искусственных 


нейронов, соединенных между собой взвешенными связями. Искусственные нейронные сети не 


программируются обычным способом согласно алгоритму, разработанному 


человеком. Они обучаются на большом числе учебных примеров с~их 


решениями (ответами) и~настраивают весовые коэффициенты (веса) связей так, 


чтобы ИНС для каждого обучающего примера давала правильный ответ. При 


этом объем матрицы весов связей может составлять несколько гигабайт. После 


обучения ИНС используется для решения задач, не входящих в~обучающие 


примеры. Что произойдет, если ИНС даст неправильный результат при их 


решении и~возникнет необходимость узнать, почему произошла ошибка, 


является ли она постоянной или нестабильной? В~традиционной программе 


человек ищет сегмент кода, ответственный за ошибку, анализирует ее 


и~исправляет. В~ИНС нет алгоритмических шагов, и~можно увидеть только 


гигабайтную матрицу весов. Эта матрица весов непрозрачна~--- с~ее помощью 


нельзя выяснить, как веса соотносятся с~ошибкой. Вот почему анализ и~оценка 


результатов, получаемых при обучении ИНС, становятся в~ряд самых 


актуальных областей исследований в~информатике~[3] и~компьютерной 


лингвистике.


   


   В статье обучение ИНС демонстрируется на материале рус\-ско-фран\-цуз\-ских 


аннотаций (см.\ примеры с~коннекторами\footnote{Коннектор~--- 


языковая единица, функция которой состоит в~выражении  


логико-се\-ман\-ти\-че\-ско\-го отношения, существующего между 


соединенными с~ее помощью частями текстового фрагмента~[4,  с.~17].} 


\textit{как (расстояние), так~и} в~табл.~\ref{t4-z} и~\textit{не только 


(расстояние)}~$\emptyset$ в~табл.~\ref{t5-z}). Аннотации были созданы во 


время эксперимента с~помощью НБД 


коннекторов~[5--10]. Каждая аннотация включает текстовый фрагмент 


с~коннектором на русском языке, перевод фрагмента на французский, 


выполненный с~помощью системы НМП Google, и~структурированное 


описание ошибок в~случае их наличия.


   


   Для эксперимента были отобраны 250 текстовых фрагментов на русском 


языке, каждый из которых содержит двухкомпонентный коннектор (см., 


например, \textit{как (расстояние), 


   так~и}~в~табл.~\ref{t4-z}). В~эксперименте использовались тексты на 


русском языке и~их французские переводы, полученные из французского 


подкорпуса Национального корпуса русского языка (НКРЯ)~[11].


   


   На протяжении года, с~марта 2019~г.\ по февраль 2020~г., для отобранных 


фрагментов в~НБД коннекторов ежемесячно формировались 


и~регистрировались 250~аннотаций. В~результате был получен массив 


экспериментальных данных, в~котором для каждого из 250~текстовых 


фрагментов была зафиксирована серия из 12~версий машинного перевода. Таким образом, всего 


в~НБД было создано 3000~аннотаций.


   


   Для обработки сформированных аннотаций применена оценка результатов 


обучения ИНС с~использованием системы количественных показателей 


(индикаторов), значения которых вычисляются на основе мониторинга работы 


ИНС (далее~--- индикаторная оценка результатов обучения). Система 


индикаторов и~методика мониторинга были разработаны в~три этапа. На первом 


этапе, описанном в~\cite{12-z}, была проведена категоризация нестабильности 


НМП на основе анализа ошибок перевода, которые возникали или исчезали 


в~сериях машинного перевода при обучении ИНС. В~результате были выделены 6~категорий~--- 


5 из них характеризуют нестабильность качества НМП:


   \begin{itemize}


   \item повышение качества НМП;


   \item снижение качества НМП;


   \item колебание качества НМП;


   \item изменение набора ошибок в~НМП без динамики его качества;


   \item изменение НМП без динамики его качества.


   \end{itemize}


   


   Шестая категория объединяет переводы, которые не менялись в~пределах 


серии (подробное описание категорий дано в~\cite{12-z}).


   


   На втором этапе была разработана и~описана методика темпоральной 


оценки нестабильности НМП~\cite{13-z}. Настоящая статья посвящена 


третьему этапу исследования, цель которого~--- оценка результатов обучения 


ИНС на основе мониторинга ее работы с~применением разработанной системы 


индикаторов (табл.~1).


   


\section{Нестабильность нейронного машинного перевода и~виды ошибок перевода}





   В результате эксперимента для каждой из 6 категорий были получены 


числовые данные для случаев, когда ошибки в~машинном переводе отсутствуют (см.\ строку~1 


табл.~\ref{t1-z}), а~также данные о~частотности 19~видов ошибок 


машинного перевода\footnote{Подробнее о~классификации ошибок см.\ в~\cite{12-z, 13-z}.} 


(см.\ строки~2--20 табл.~\ref{t1-z}). Частотность вычислялась от общего числа 


ошибок (2836), зарегистрированных в~1800 аннотациях, так как  одна 


аннотация могла содержать описание нескольких ошибок.





  На момент написания настоящей статьи была проведена рубрикация по 


видам ошибок и~категориям, характеризующим нестабильность НМП, для 


60\% сформированных в~НБД аннотаций (1800 из~3000). На основании анализа 


полученных данных можно выявить ряд закономерностей. 


  


  По сумме ошибок во всех 6 категориях, характеризующих 


нестабильность НМП, самой частотной стала ошибка ErrorSemant~--- 


лексическая ошибка (38,26\% от 2836~ошибок), искажающая смысл оригинала 


в~переводном фрагменте (см.\ также строку~11 в~табл.~\ref{t4-z} и~строки~6 


и~11 в~табл.~\ref{t5-z}). В~НБД было зарегистрировано 1085~аннотаций 


с~такой ошибкой из~1800.





\pagebreak


  


\begin{table}\footnotesize %tabl1


\begin{center}


\Caption{Индикаторы нестабильности НМП, ошибок и~случаев их отсутствия}


\label{t1-z}


\vspace*{2ex}





\tabcolsep=1.55pt 


\begin{tabular}{|r|l|c|c|c|c|c|c|c|c|}


\hline


№&\multicolumn{1}{c|}{\tabcolsep=0pt 


\begin{tabular}{c} 


Код ошибки\end{tabular}}&\tabcolsep=0pt \begin{tabular}{c}Повы-\\ шение \\ 


каче-\\ ства\\  НМП \end{tabular}&\tabcolsep=0pt \begin{tabular}{c}Сни-\\ жение \\ 


каче-\\ ства \\ НМП \end{tabular}&\tabcolsep=0pt \begin{tabular}{c}Коле-\\ бание \\ 


каче-\\ ства \\ НМП \end{tabular}&\tabcolsep=0pt \begin{tabular}{c} Измене-\\ ние\\ 


набора\\ ошибок\\ в~НМП\\ без дина-\\ мики его \\ качества 


\end{tabular}&\tabcolsep=0pt \begin{tabular}{c}Изме-\\ нение\\ НМП\\ без дина-\\ 


мики его \\ качества \end{tabular}&\tabcolsep=0pt \begin{tabular}{c}НМП\\ без\\ 


изме-\\ нений\end{tabular}&\tabcolsep=0pt \begin{tabular}{c} 


{Сумма}\\  {по}\\ {6 кате-}\\ { гориям,}\\ 


{ характери-}\\ { зующим}\\ { нестабиль-}\\ { ность}\end{tabular}&


\tabcolsep=0pt \begin{tabular}{c} 


Значе-\\ ния\\  индика-\\ торов\\  ошибок,\\ \%\end{tabular}\\


\cline{3-8}


&&1&2&3&4&5&6&&\\  


\hline


1&NoError&73&66\hphantom{9}&82&\hphantom{9}0&156\hphantom{9}&48\hphantom{9}


&425\hphantom{9}&14,99\hphantom{9}\\


2&ErrorSemant&366\hphantom{9}&67\hphantom{9}&362\hphantom{9}&217\hphantom{9}


&61&12\hphantom{9}&1085\hphantom{99}&38,26\hphantom{9}\\


3&ErrorSyntax&108\hphantom{9}&25\hphantom{9}&94&66&12&0&305\hphantom{9}&10,75\hphantom{9}\\


4&ErrorSyntaxPostCNT&82&19\hphantom{9}&60&46&24&0&231\hphantom{9}&8,15\\


5&ErrorTotalCNT&64&3&19&62&24&0&172\hphantom{9}&6,06\\


6&ErrorCNT&46&0&33&55&\hphantom{9}0&0&134\hphantom{9}&4,72\\


7&Pleonasm&28&12\hphantom{9}&15&23&11&0&89&3,14\\


8&ErrorMorphPostCNT&16&0&17&27&\hphantom{9}0&0&60&2,12\\


9&ErrorMorphCNT&16&0&\hphantom{9}2&20&12&0&50&1,76\\


10&Lacuna&13&0&16&36&\hphantom{9}0&0&65&2,29\\


11&ErrorPart2CNT&26&3&\hphantom{9}6&16&12&0&63&2,22\\


12&ErrorPunct&\hphantom{9}9&3&21&10&\hphantom{9}0&0&43&1,52\\


13&ErrorPart1CNT&15&0&16&13&12&0&56&1,97\\


14&ErrorMorph&\hphantom{9}5&3&14&13&\hphantom{9}0&0&35&1,23\\


15&AgramTotal&\hphantom{9}5&0&\hphantom{9}0&\hphantom{9}7&\hphantom{9}0&0&12&0,42\\


16&TrPart1CNT&\hphantom{9}0&0&\hphantom{9}5&\hphantom{9}0&\hphantom{9}0&0&\hphantom{9}5&0,18\\


17&Latin&\hphantom{9}5&0&\hphantom{9}0&\hphantom{9}0&\hphantom{9}0&


0&\hphantom{9}5&0,18\\


18&TrPart2CNT&\hphantom{9}0&0&\hphantom{9}1&\hphantom{9}0&\hphantom{9}0&0&\hphantom{9}1&0,04\\


19&Cyrillic&\hphantom{9}0&0&\hphantom{9}0&\hphantom{9}0&\hphantom{9}0


&0&\hphantom{9}0&0,00\\


20&ErrorOrthCNT&\hphantom{9}0&0&\hphantom{9}0&\hphantom{9}0&\hphantom{9}0&0&\hphantom{9}0&0,00\\


\hline


\multicolumn{2}{|l|}{\tabcolsep=0pt\begin{tabular}{l} 


Сумма по видам\\ 


 ошибок\end{tabular}}& 


877\hphantom{9}&201\hphantom{99}&763\hphantom{9}&611\hphantom{9}&


324\hphantom{9}&60\hphantom{9}&2836\hphantom{99}&\\


\hline


\multicolumn{2}{|l|}{\tabcolsep=0pt\begin{tabular}{l} 


Значения индикаторов\\ 


 для каждой категории, 


\%\end{tabular}}&30,92&7,09&26,90&21,54&11,42&2,12&&100,00\\


\hline 


   \end{tabular}


   \end{center}


   \end{table}


   


  Далее по уменьшению частотности (пропуская код NoError, 


ха\-рак\-те\-ри\-зу\-ющий допустимые варианты перевода)~--- синтаксические ошибки 


во фрагменте, не вводимом коннектором (ErrorSyntax~--- 10,75\%) и~вводимом 


коннектором (\mbox{ErrorSyntaxPostCNT}~--- 8,15\%). Пример перевода с~ошибкой 


\mbox{ErrorSyntaxPostCNT} указан в~строке~1 табл.~\ref{t4-z}. Морфологических 


ошибок (\mbox{ErrorMorphPostCNT} для фрагмента с~коннектором и~\mbox{ErrorMorph} для 


фрагмента, не вводимого коннектором), напротив, значительно меньше~--- 


2,12\% и~1,76\% соответственно. Это свидетельствует о~том, что система НМП 


в~основном хорошо справляется с~грамматикой на уровне морфологии 


и~трудности возникают чаще на уровне синтаксиса. Однако устойчивое 


снижение числа морфологических ошибок в~период проведения эксперимента 


не наблюдалось (табл.~\ref{t2-z}).


  


  Среди ошибок, касающихся непосредственно перевода коннекторов, самыми 


распространенными стали ErrorTotalCNT (система НМП переводит коннектор 


несуществующей языковой единицей, пример см.\ в~строке~1 табл.~\ref{t4-z}) 


и~ErrorCNT (система определяет, что в~переводе нужен коннектор, но 


ошибается в~его выборе, см.\ строку~6 табл.~\ref{t4-z}).


  


\begin{table}\small %tabl2


\begin{center}


\Caption{Динамика ошибок по месяцам}


\label{t2-z}


\vspace*{2ex}





\tabcolsep=3.5pt 


\begin{tabular}{|l|c|c|c|c|c|c|c|c|c|c|c|c|c|}


\hline


\multicolumn{1}{|c|}{Код}&\multicolumn{12}{c|}{ Месяц}\\ 


\cline{2-13}


\multicolumn{1}{|c|}{ошибки}&1&2&3&4&5&6&7&8&9&10&11&12\\


\hline


AgramTotal&3&2&3&2&2&0&0&0&0&0&0&0\\


ErrorMorph&3&2&3&2&2&2&2&3&3&5&4&4\\


ErrorMorphPostCNT&5&5&7&6&7&4&4&4&5&5&4&4\\


ErrorSyntax&28\hphantom{9}&27\hphantom{9}&28\hphantom{9}&26\hphantom{9}


&27\hphantom{9}&23\hphantom{9}&22\hphantom{9}&21\hphantom{9}&


24\hphantom{9}&28\hphantom{9}&25\hphantom{9}&26\hphantom{9}\\


ErrorSyntaxPostCNT&27\hphantom{9}&24\hphantom{9}&22\hphantom{9}&


25\hphantom{9}&22\hphantom{9}&15\hphantom{9}&16\hphantom{9}&16\hphantom{9}


&16\hphantom{9}&16\hphantom{9}&16\hphantom{9}&16\hphantom{9}\\


ErrorSemant&85\hphantom{9}&84\hphantom{9}&82\hphantom{9}&83\hphantom{9}


&83\hphantom{9}&93\hphantom{9}&94\hphantom{9}&95\hphantom{9}&96\hphantom{9}


&97\hphantom{9}&95\hphantom{9}&98\hphantom{9}\\


ErrorPunct&4&5&4&4&4&3&3&3&3&2&4&4\\


Cyrillic&0&0&0&0&0&0&0&0&0&0&0&0\\


Latin&1&1&1&1&1&0&0&0&0&0&0&0\\


Lacuna&8&5&8&5&4&6&6&6&6&5&3&3\\


Pleonasm&6&7&9&10\hphantom{9}&8&7&8&8&7&6&6&7\\


TrPart1CNT&0&0&0&0&0&2&1&1&1&0&0&0\\


TrPart2CNT&0&0&0&1&0&0&0&0&0&1&1&1\\


ErrorPart1CNT&5&5&6&6&6&2&2&2&3&4&5&5\\


ErrorPart2CNT&10\hphantom{9}&10\hphantom{9}&8&8&8&2&2&2&2&2&2&2\\


ErrorTotalCNT&20\hphantom{9}&21\hphantom{9}&22\hphantom{9}&22\hphantom{9}


&23\hphantom{9}&7&7&8&9&8&9&9\\


ErrorMorphCNT&6&5&5&6&5&3&3&2&2&3&4&6\\


ErrorOrthCNT&0&0&0&0&0&0&0&0&0&0&0&0\\


ErrorCNT&13\hphantom{9}&13\hphantom{9}&14\hphantom{9}&14\hphantom{9}


&15\hphantom{9}&9&9&9&11\hphantom{9}&9&9&9\\


NoError&32\hphantom{9}&34\hphantom{9}&36\hphantom{9}&37\hphantom{9}&


35\hphantom{9}&39\hphantom{9}&39\hphantom{9}&38\hphantom{9}&36\hphantom{9}


&32\hphantom{9}&34\hphantom{9}&33\hphantom{9}\\


\hline


\end{tabular}


\end{center}


\end{table}





  Хотя система НМП нередко допускает ошибки при переводе коннектора на 


французский язык, по данным из НБД коннекторов, не вошедшим 


в~табл.~\ref{t1-z}, было зарегистрировано относительно мало случаев, когда 


коннектор не переведен вовсе, всего 34~аннотации из~1800. Объем настоящей 


статьи не позволяет описать это наблюдение подробнее, отметим только, что 


перевод русскоязычного коннектора нулевым эквивалентом не всегда 


ошибочен.


  


  На данном этапе нулевыми остались ячейки для ошибок Cyrillic (слова 


кириллицей в~переводе) и~ErrorOrthCNT (орфографическая ошибка в~форме 


коннектора). Это может указывать на тенденцию к~повышению качества НМП, 


поскольку в~предыдущих исследованиях~\cite{14-z} эти значения не были 


нулевыми.


  


  Кроме того, рассматривая данные для первых пяти категорий, которые 


характеризуют нестабильность НМП, можно отметить следующее: первой 


категории <<Повышение качества НМП>> соответствует наибольшее число 


ошибок (30,92\%), а~второй категории <<Снижение качества НМП>>~--- 


наименьшее число ошибок (7,09\%), что заслуживает отдельного 


лингвистического анализа, выходящего за рамки этой статьи.


  


\section{Динамика ошибок на протяжении эксперимента}





  Использование НБД при изучении нестабильности НМП позволило не только 


получить значения индикаторов по категориям, характеризующим 


нестабильность НМП, и~по ошибкам НМП, но и~проследить динамику ошибок 


каждого вида на протяжении всего эксперимента (см.\ табл.~\ref{t2-z}).


  


  Данные табл.~\ref{t2-z} свидетельствуют о~том, что у разных видов ошибок 


разная динамика. Основная тенденция (строки для ошибок ErrorSyntaxPostCNT, 


Lacuna, ErrorPart2CNT, ErrorTotalCNT 


и~ErrorCNT\footnote{ErrorSyntaxPostCNT~--- cинтаксическая ошибка во 


фрагменте текста, вводимом коннектором; Lacuna~--- пропуск фрагмента 


текста; ErrorPart2CNT~--- вторая часть неоднословного коннектора 


переведена ошибочно; ErrorTotalCNT~--- коннектор переведен 


несуществующей языковой единицей; ErrorCNT~--- семантическая ошибка 


в~выборе коннектора.})~--- уменьшение числа ошибок со временем. Однако 


есть ошибки (ErrorSyntax, Pleonasm\footnote{ErrorSyntax~--- синтаксическая 


ошибка во фрагменте текста, не вводимом коннектором; Pleonasm~---  


избыточный перевод.} и~\mbox{ErrorMorphPostCNT}, 


ErrorMorph\footnote{ErrorMorphPostCNT~--- морфологическая ошибка во 


фрагменте текста, вводимом коннектором; ErrorMorph~--- морфологическая 


ошибка во фрагменте текста, не вводимом коннектором.}), число которых 


колеблется. Также есть виды ошибок, число которых, напротив, увеличивается 


(например, ErrorSemant\footnote{ErrorSemant~--- лексическая ошибка, 


искажение смысла переводимого фрагмента.} на 15,29\%).





\section{Распределение аннотаций по категориям, характеризующим 


нестабильность нейронного машинного перевода}





  Тенденции, продемонстрированные в~табл.~\ref{t2-z}, коррелируют 


с~данными табл.~\ref{t3-z}, в~которой представлено распределение аннотаций 


по категориям, ха\-рак\-те\-ри\-зу\-ющим нестабильность НМП.


  


  Так, наибольшее значение индикатора (30,6\% аннотаций) у 1-й категории, 


характеризующей повышение качества НМП (пример см.\ в~табл.~\ref{t4-z}). 


За ней по частотности следуют рубрики, сигнализирующие об изменении НМП 


без однозначной динамики его качества к~повышению или понижению. Это 


колебание качества НМП (3-я категория), изменение набора ошибок в~НМП 


без динамики его качества (4-я категория) и~изменение НМП без динамики его 


качества (5-я категория). Наименьшее число аннотаций приходится на 2-ю 


категорию (снижение качества НМП)~--- 9,3\% аннотаций (пример см.\ 


в~табл.~\ref{t5-z}). Стабильными в~течение эксперимента остались лишь 3,3\% 


переводов (НМП без изменений).


  





\begin{table}\small %tabl3


\begin{center}


\Caption{Распределение аннотаций по категориям, характеризующим 


нестабильность НМП


\label{t3-z}}


\vspace*{3pt}





\tabcolsep=7.9pt 


\begin{tabular}{|c|l|c|c|}


\hline


№&\multicolumn{1}{c|}{\tabcolsep=0pt\begin{tabular}{c}Название рубрики,


характеризующей\\ нестабильность НМП\end{tabular}}&


\tabcolsep=0pt\begin{tabular}{c}Число\\ аннотаций\end{tabular}&


\tabcolsep=0pt\begin{tabular}{c}Доля от общего числа\\ аннотаций (1800) 


\end{tabular}\\


\hline


1&Повышение качества НМП&552&30,6\%\\


\hline


2&Снижение качества НМП&168&\hphantom{9}9,3\%\\


\hline


3&Колебание качества НМП&492&27,3\%\\


\hline


4&\tabcolsep=0pt\begin{tabular}{l}Изменение набора ошибок в~НМП\\ без 


динамики его качества\end{tabular}& 300&16,6\%\\


\hline


5&\tabcolsep=0pt\begin{tabular}{l}Изменение НМП без динамики\\ его 


качества\end{tabular}& 228&12,6\%\\


\hline


6&НМП без изменений&\hphantom{9}60&\hphantom{9}3,3\%\\


\hline


\end{tabular}


\end{center}


\end{table}


  


  В табл.~\ref{t4-z} представлена серия переводов следующего текстового 


фрагмента: <<Как воровали, так и~будут воровать>> [Светлана Алексиевич. 


Время секонд хэнд (ч.~2) (2013)]\footnote[1]{Профессиональный перевод этого 


предложения на французский язык: <<Les gens continueront 


$\grave{\mbox{a}}$~voler comme avant>> [Sophie Benech, 2013].}. По итогам 


проведенного лингвистического анализа, методика которого  подробно описана


в~\cite{13-z}, для этой серии переводов была проставлена категория 


<<Повышение качества НМП>>.


  


\begin{table}\footnotesize %tabl4


\vspace*{-4pt}


\begin{center}


\Caption{Серия аннотаций с~повышением качества НМП


\label{t4-z}}


\vspace*{2ex}





\tabcolsep=2.4pt 


\begin{tabular}{|c|c|l|l|l|l|c|c|}


\hline


\tabcolsep=0pt \begin{tabular}{c}№\\ п/п\end{tabular}&


\tabcolsep=0pt \begin{tabular}{c}№\\ в~НБД\end{tabular}&


\multicolumn{1}{c|}{\tabcolsep=0pt \begin{tabular}{c} 


Контекст РР \\ в~оригинале$^1$\end{tabular}} &


\multicolumn{1}{c|}{\tabcolsep=0pt \begin{tabular}{c} 


РР в~ори-\\ гинале\\  и~ее при-\\ знаки \end{tabular}}&


\multicolumn{1}{c|}{\tabcolsep=0pt \begin{tabular}{c} 


 Контекст РР\\ в~переводе$^2$\end{tabular} }&


\multicolumn{1}{c|}{\tabcolsep=0pt \begin{tabular}{c}


 РР в~переводе\\ и ее признаки\end{tabular}}&


\multicolumn{1}{c|}{\tabcolsep=0pt \begin{tabular}{c} 


При-\\ знаки \\  анно-\\ тации\end{tabular}}&


\multicolumn{1}{c|}{\tabcolsep=0pt \begin{tabular}{c} 


Дата\\ и~время\\  записи\\ в~НБД\end{tabular}}\\ 


\hline


\hphantom{1}1&28312&\tabcolsep=0pt\begin{tabular}{l} 


\textbf{Как}\\ воровали,\\ \textbf{так и}~будут\\ воровать. \end{tabular}&


\tabcolsep=0pt \begin{tabular}{l} 


\textbf{как}$\|$\textbf{так и}\\ 


$\langle$аналогия$\rangle$\\ 


$\langle$CNT p\\  


CNT q$\rangle$


%\\ $\langle$CNT$\rangle$\\ 


%$\langle$Дистант$\rangle$ 


\end{tabular} &


\tabcolsep=0pt\begin{tabular}{l} 


\textbf{Comment} \\ voler, \textbf{alors}\\ ils voleront. \end{tabular}&


\tabcolsep=0pt \begin{tabular}{l} 


\textbf{comment}$\|$\textbf{alors}\\  


$\langle$ErrorSyntaxPostCNT$\rangle$\\ 


$\langle$ErrorTotalCNT$\rangle$\\ 


$\langle$ErrorMorphCNT$\rangle$ 


\end{tabular} 


&MT&


\tabcolsep=0pt\begin{tabular}{c} 03.03.19\\ 17:20\end{tabular}\\ 


\cline{1-2} 


\cline{5-8}


\hphantom{1}2&28670&&


\tabcolsep=0pt \begin{tabular}{l} 


$\langle$CNT$\rangle$\\ 


$\langle$Дистант$\rangle$ 


\end{tabular} 


&\multicolumn{1}{c|}{-/\!\!/-}&\multicolumn{1}{c|}{-/\!\!/-}&-/\!\!/-&


\tabcolsep=0pt\begin{tabular}{c} 01.04.19\\ 15:39\end{tabular}\\ 


\cline{1-2} 


\cline{5-8}


\hphantom{1}3&29324&&&\multicolumn{1}{c|}{-/\!\!/-}&\multicolumn{1}{c|}{-/\!\!/-}&-/\!\!/-&


\tabcolsep=0pt\begin{tabular}{c} 05.05.19\\ 14:38\end{tabular}\\ 


\cline{1-2} 


\cline{5-8}


\hphantom{1}4&29913&&&\multicolumn{1}{c|}{-/\!\!/-}&\multicolumn{1}{c|}{-/\!\!/-}&-/\!\!/-&


\tabcolsep=0pt\begin{tabular}{c} 05.06.19\\ 14:26\end{tabular}\\ 


\cline{1-2} 


\cline{5-8}


\hphantom{1}5&30174&&&\multicolumn{1}{c|}{-/\!\!/-}&\multicolumn{1}{c|}{-/\!\!/-}&-/\!\!/-&


\tabcolsep=0pt\begin{tabular}{c} 02.07.19\\ 13:45\end{tabular}\\ 


\cline{1-2} 


\cline{5-8}


\hphantom{1}6&30461&&&\tabcolsep=0pt\begin{tabular}{l} 


\textbf{Comme} ils\\ ont vol$\acute{\mbox{e}}$,\\ ils vont voler. 


\end{tabular}&


\tabcolsep=0pt \begin{tabular}{l} 


\textbf{comme}\\  


$\langle$ErrorCNT$\rangle$ 


\end{tabular}&


\tabcolsep=0pt \begin{tabular}{c} Cngrn\\ MT\end{tabular}&


\tabcolsep=0pt\begin{tabular}{c} 03.08.19\\ 23:37\end{tabular}\\ 


\cline{1-2} 


\cline{5-8}


\hphantom{1}7&30707&&&\multicolumn{1}{c|}{-/\!\!/-}&\multicolumn{1}{c|}{-/\!\!/-}&-/\!\!/-&


\tabcolsep=0pt\begin{tabular}{c} 03.09.19\\ 15:04\end{tabular}\\ 


\cline{1-2} 


\cline{5-8}


\hphantom{1}8&31022&&&\multicolumn{1}{c|}{-/\!\!/-}&\multicolumn{1}{c|}{-/\!\!/-}&-/\!\!/-&


\tabcolsep=0pt\begin{tabular}{c} 02.10.19\\ 22:23\end{tabular}\\ 


\cline{1-2} 


\cline{5-8}


\hphantom{1}9&31545&&&\multicolumn{1}{c|}{-/\!\!/-}&\multicolumn{1}{c|}{-/\!\!/-}&-/\!\!/-&


\tabcolsep=0pt\begin{tabular}{c} 04.11.19\\ 15:27\end{tabular}\\ 


\cline{1-2} 


\cline{5-8}


10&32078&&&\multicolumn{1}{c|}{-/\!\!/-}&\multicolumn{1}{c|}{-/\!\!/-}&-/\!\!/-&


\tabcolsep=0pt\begin{tabular}{c} 02.12.19\\ 15:36\end{tabular}\\ 


\cline{1-2} 


\cline{5-8}


11&32485&&&\tabcolsep=0pt \begin{tabular}{l}\textbf{En volant},\\ ils voleront. \end{tabular}&


\multicolumn{1}{l|}{\tabcolsep=0pt \begin{tabular}{l} 


\textbf{x\_g$\acute{\mbox{e}}$r\_pr$\acute{\mbox{e}}$s}\\ 


$\langle$ErrorSemant$\rangle$ \end{tabular}}&


\tabcolsep=0pt \begin{tabular}{c} Dvrg\\ MT \end{tabular}&


\tabcolsep=0pt\begin{tabular}{c} 02.01.20\\ 20:02\end{tabular}\\ 


\cline{1-2} 


\cline{5-8}


12&33171&&&\multicolumn{1}{c|}{-/\!\!/-}&


\multicolumn{1}{l|}{\tabcolsep=0pt \begin{tabular}{l} 


\textbf{x\_g$\acute{\mbox{e}}$r\_pr$\acute{\mbox{e}}$s}\\  


$\langle$ErrorSemant$\rangle$\\ 


$\langle$Повышение\\ качества НМП$\rangle$ \end{tabular}}&


-/\!\!/-&\tabcolsep=0pt\begin{tabular}{c} 03.02.20\\ 15:50\end{tabular}\\ 


\hline


  \multicolumn{8}{p{367.9pt}}{\scriptsize


  \hspace*{2mm}$^1$Речевая реализация (РР) в~оригинале~--- форма коннектора 


в~конкретном высказывании~\cite{15-z}.}\\


  \multicolumn{8}{p{367.9pt}}{\scriptsize %\footnotesize


  \hspace*{2mm}$^2$Речевая реализация в~переводе~--- функционально эквивалентный 


фрагмент перевода, передающий смысловое наполнение РР в~оригинале. 


Понятие <<функционально эквивалентный фрагмент>> предложено в~\cite{16-z}.}\\


  %\multicolumn{8}{p{367.9pt}}{\footnotesize


 % $^3$Подробнее о~рубриках см.\ в~\cite{6-z}.}\\


  \multicolumn{8}{p{367.9pt}}{\scriptsize \textbf{Примечания.}\vspace*{-3pt}


  \begin{enumerate}[1.]


\item Рубрика (подробнее о~рубриках см.\ в~\cite{6-z}) 


<<аналогия>> говорит о~том, что коннектор \textit{как (расстояние), так 


и} выражает ло\-ги\-ко-се\-ман\-ти\-че\-ское отношение аналогии.


\vspace*{-3pt}


\item Рубрика <<CNT p CNT q>> говорит о~том, что элементы 


двухкомпонентного коннектора находятся в~каждом из соединяемых 


компонентов текста~$p$ (<<воровали>>) и~$q$ (<<будут воровать>>).\vspace*{-3pt}


\item Рубрика <<CNT>> говорит о~том, что аннотация сформирована для 


всего коннектора, а~не для отдельно составляющих его блоков или 


элементов.\vspace*{-3pt}


\item Рубрика <<Дистант>> говорит о~том, что составные части 


коннектора <<как>> и~<<так и>> разделены текстом <<воровали>>.\vspace*{-3pt}


\item Рубрика <<MT>> проставляется автоматически при использовании 


машинного переводчика.\vspace*{-3pt}


\item Рубрика <<Cngrn>> говорит о~том, что коннектор русского языка 


переведен именно коннектором французского языка, а~не языковой 


единицей или конструкцией другой категории.\vspace*{-3pt}


\item Рубрика <<Dvrg>> говорит о~том, что коннектор русского языка 


переведен языковой единицей или конструкцией другой категории (не 


коннектором).\vspace*{-3pt}


\item Рубрика <<Повышение качества НМП>> и~другие рубрики, характеризующие 


изменение качества перевода, проставляются в~последней, 12-й версии 


перевода, но характеризуют всю серию из 12 переводов.


\end{enumerate}}


\end{tabular}


\vspace*{-14.1pt}


\end{center}


\vspace*{-3pt}


\end{table}





  За год эксперимента по изучению нестабильности НМП, которая, по всей 


видимости, вызвана обучением ИНС, машинный перевод исходного текстового фрагмента 


менялся три раза (повторяющиеся значения таблицы заменены символом -/\!\!/-). 


В~первой версии перевода от 03.03.19 допущено три ошибки. Это 


\mbox{ErrorSyntaxPostCNT} (отсутствует синтаксическая связь между клаузами 


\textit{comment voler} и~\textit{alors ils voleront}), а~также две ошибки 


в~переводе коннектора~--- ErrorTotalCNT (русский коннектор 


\textit{как}$\|$\textit{так~и} переведен несуществующей во французском языке 


единицей \textit{comment}$\|$\textit{alors}) и~ErrorMorph (часть коннектора 


ошибочно заменена языковой единицей \textit{comment}, не являющейся коннекто-\linebreak


ром). 


  


  В 6-м переводе от 03.08.19 допущена всего одна ошибка~--- это 


ErrorCNT (выбран семантически неверный коннектор \textit{comme}). 


Подобные ошибки зачастую ведут к~искажению смысла всего переводимого 


фрагмента: при использовании коннектора \textit{comme} фраза приобретает 


смысл <<Потому как они воровали, они будут воровать>>. 


  


  В 11-м переводе от 02.01.20 допущена ErrorSemant: при переводе 


коннектора \textit{как}$\|$\textit{так~и} используется грамматическая 


структура g$\acute{\mbox{e}}$rondif pr$\acute{\mbox{e}}$sent, и~фраза 


приобретает смысл <<Воруя, они будут воровать>>.


  


  Несмотря на то что в~последнем, 12-м, переводе сохраняется 


ErrorSemant, в~целом в~этой серии переводов прослеживается тенденция 


к~повышению качества мащинного перевода, так как число ошибок со временем уменьшается 


и~переводной текст становится более понятным. 


  


\begin{table}\footnotesize %tabl5


\begin{center}


\Caption{Серия аннотаций со снижением качества НМП


\label{t5-z}}


\vspace*{2ex} 





\tabcolsep=1.35pt 


\begin{tabular}{|c|c|l|l|l|l|c|l|}


\hline


\tabcolsep=0pt \begin{tabular}{c}№\\ п/п\end{tabular}&


\tabcolsep=0pt \begin{tabular}{c}№\\ в~НБД\end{tabular}&


\multicolumn{1}{c|}{\tabcolsep=0pt \begin{tabular}{c} 


Контекст РР \\ в~оригинале\end{tabular}} &


\multicolumn{1}{c|}{\tabcolsep=0pt \begin{tabular}{c} 


РР в~ори-\\ гинале\\  и~ее при-\\ знаки \end{tabular}}&


\multicolumn{1}{c|}{\tabcolsep=0pt \begin{tabular}{c} 


 Контекст РР\\ в~переводе\end{tabular} }&


\multicolumn{1}{c|}{\tabcolsep=0pt \begin{tabular}{c}


 РР в~переводе\\ и ее признаки\end{tabular}}&


\multicolumn{1}{c|}{\tabcolsep=0pt \begin{tabular}{c} 


При-\\ знаки \\  анно-\\ тации\end{tabular}}&


\multicolumn{1}{c|}{\tabcolsep=0pt \begin{tabular}{c} 


Дата\\ и~время\\  записи\\ в~НБД\end{tabular}}\\ 


\hline


&&&&&&&\\[-9pt]


\hphantom{1}1&28213&\tabcolsep=0pt\begin{tabular}{l} 


Каримов сказал:\\ 


---~У~нас уничто-\\ жили 


\textbf{не только} \\людей, нацио-\\ нальную культуру\\ 


уничтожили. \end{tabular}


&


\tabcolsep=0pt \begin{tabular}{l} 


\textbf{не только}$\|$$\boldsymbol{\emptyset}$\\ 


$\langle$неединст-\\ венности$\rangle$\\ 


$\langle$CNT p\\  


CNT q$\rangle$\\ 


$\langle$CNT$\rangle$


%\\ $\langle$Контакт$\rangle$ 


\end{tabular} 


&\tabcolsep=0pt\begin{tabular}{l} 


Karimov a d$\acute{\mbox{e}}$clar$\acute{\mbox{e}}$:\\ 


<<\textbf{Pas seulement}\\ les gens ont $\acute{\mbox{e}}$t$\acute{\mbox{e}}$\\ 


d$\acute{\mbox{e}}$truits ici,\\ la culture nationale\\ 


a $\acute{\mbox{e}}$t$\acute{\mbox{e}}$ 


 d$\acute{\mbox{e}}$truite>>.\end{tabular}&


\tabcolsep=0pt\begin{tabular}{l}\textbf{pas}\\ \textbf{seulement}$\|$$\boldsymbol{\emptyset}$ \\ 


$\langle$NoError $\rangle$ 


\end{tabular} 


&\tabcolsep=0pt 


\begin{tabular}{c}Cngrn\\ MT\end{tabular}&


\tabcolsep=0pt \begin{tabular}{c} 01.03.19\\ 15:18 \end{tabular}\\ 


\cline{1-2} 


\cline{5-8}


\hphantom{1}2&28724&&


\tabcolsep=0pt \begin{tabular}{l} 


$\langle$Контакт$\rangle$\\ \ 


\end{tabular} 


&\multicolumn{1}{c|}{-/\!\!/-}&\multicolumn{1}{c|}{-/\!\!/-}&-/\!\!/-&


\tabcolsep=0pt \begin{tabular}{c} 02.04.19\\ 14:49\end{tabular}\\ 


\cline{1-2} 


\cline{5-8}


\hphantom{1}3&29413&&&\multicolumn{1}{c|}{-/\!\!/-}&\multicolumn{1}{c|}{-/\!\!/-}&-/\!\!/-&


\tabcolsep=0pt \begin{tabular}{c} 06.05.19\\ 14:45\end{tabular}\\ 


\cline{1-2} 


\cline{5-8}


\hphantom{1}4&30049&&&\multicolumn{1}{c|}{-/\!\!/-}&\multicolumn{1}{c|}{-/\!\!/-}&-/\!\!/-&


\tabcolsep=0pt \begin{tabular}{c} 07.06.19\\ 14:31\end{tabular}\\ 


\cline{1-2} 


\cline{5-8}


\hphantom{1}5&30340&&&\multicolumn{1}{c|}{-/\!\!/-}&\multicolumn{1}{c|}{-/\!\!/-}&-/\!\!/-&


\tabcolsep=0pt \begin{tabular}{c} 05.07.19\\ 14:28\end{tabular}\\ 


\cline{1-2} 


\cline{5-8}


&&&&&&&\\[-9pt]


\hphantom{1}6&30557&&&\tabcolsep=0pt\begin{tabular}{l} 


Karimov a d$\acute{\mbox{e}}$clar$\acute{\mbox{e}}$: \\ 


---~Nous \textbf{n'}avons\\ \textbf{pas seulement}\\ d$\acute{\mbox{e}}$truit des\\ 


personnes, nous\\ avons $\acute{\mbox{\textbf{e}}}$\textbf{galement}\\ 


 d$\acute{\mbox{e}}$truit la culture\\ nationale. \end{tabular}& 


\tabcolsep=0pt\begin{tabular}{l} 


\textbf{ne}$\vert$\textbf{pas}\\ \textbf{seulement}$\|$\\ $\acute{\mbox{\textbf{e}}}$\textbf{galement}\\


$\langle$ErrorSemant$\rangle$ 


\end{tabular} 


&-/\!\!/-&


\tabcolsep=0pt \begin{tabular}{c} 05.08.19\\ 22:29\end{tabular}\\ 


\cline{1-2} 


\cline{5-8}


\hphantom{1}7&30761&&&\multicolumn{1}{c|}{-/\!\!/-}&\multicolumn{1}{c|}{-/\!\!/-}&-/\!\!/-&


\tabcolsep=0pt \begin{tabular}{c} 04.09.19\\ 15:08\end{tabular}\\ 


\cline{1-2} 


\cline{5-8}


\hphantom{1}8&31181&&&\multicolumn{1}{c|}{-/\!\!/-}&\multicolumn{1}{c|}{-/\!\!/-}&-/\!\!/-&


\tabcolsep=0pt \begin{tabular}{c} 06.10.19\\ 23:28\end{tabular}\\ 


\cline{1-2} 


\cline{5-8}


\hphantom{1}9&31599&&&\multicolumn{1}{c|}{-/\!\!/-}&\multicolumn{1}{c|}{-/\!\!/-}&-/\!\!/-&


\tabcolsep=0pt \begin{tabular}{c} 05.11.19\\ 14:39\end{tabular}\\ 


\cline{1-2} 


\cline{5-8}


10&32186&&&\multicolumn{1}{c|}{-/\!\!/-}&\multicolumn{1}{c|}{-/\!\!/-}&-/\!\!/-&


\tabcolsep=0pt \begin{tabular}{c} 03.12.19\\ 15:08\end{tabular}\\ 


\cline{1-2} 


\cline{5-8}


&&&&&&&\\[-9pt]


11&32594&&&\tabcolsep=0pt\begin{tabular}{l} 


Karimov a d$\acute{\mbox{e}}$clar$\acute{\mbox{e}}$:\\  


---~Nous avons \textbf{non}\\ \textbf{seulement} d$\acute{\mbox{e}}$truit\\ 


 des gens, nous\\ avons d$\acute{\mbox{e}}$truit la\\ culture nationale. \end{tabular}& 


\tabcolsep=0pt\begin{tabular}{l} 


\textbf{non}\\ \textbf{seulement}$\|$$\boldsymbol{\emptyset}$\\ 


$\langle$ErrorSemant$\rangle$ 


\end{tabular} 


&-/\!\!/-&


\tabcolsep=0pt \begin{tabular}{c} 04.01.20\\ 11:51\end{tabular}\\ 


\cline{1-2} 


\cline{5-8}


12&33253&&&\multicolumn{1}{c|}{-/\!\!/-}&\tabcolsep=0pt\begin{tabular}{l} 


\textbf{non}\\ \textbf{seulement}$\|$$\boldsymbol{\emptyset}$\\  


$\langle$ErrorSemant$\rangle$\\ 


$\langle$Снижение ка-\\ чества НМП$\rangle$ 


\end{tabular} 


&-/\!\!/-&


\tabcolsep=0pt \begin{tabular}{c} 05.02.20\\ 22:40\end{tabular}\\ 


\hline


  \end{tabular}


  \end{center}


  \end{table}





\pagebreak


  


  В табл.~\ref{t5-z} представлена серия переводов следующего русскоязычного 


фрагмента: <<Каримов сказал: ---~У~нас уничтожили не только людей, 


национальную культуру уничтожили>> [В.\,С.~Гроссман. Жизнь и~судьба 


(ч.~1) (1959)]\footnote{ Профессиональный перевод этого предложения на 


французский язык: <<---~Mais chez nous, reprit Karimov, on ne s'est pas 


content$\acute{\mbox{e}}$ d'an$\acute{\mbox{e}}$antir des hommes, on a 


an$\acute{\mbox{e}}$anti toute une culture>> [Alexis Berelowitch, 1980].}. 


  


 


  За год перевод исходного фрагмента также менялся 3 раза. Первая версия 


от 01.03.19~--- приемлемый вариант перевода, поэтому для нее проставлена 


рубрика \mbox{NoError}. В~6-й версии от 05.08.19 проставлена рубрика 


\mbox{ErrorSemant}: в~переводе искажен смысл исходного фрагмента из-за ошибки 


в~выборе семантического субъекта (вместо <<у~нас уничтожили>> получилось 


<<мы уничтожили>>). В~11-й версии перевода \textit{des personnes} 


заменяется на \textit{des gens}, но указанная \mbox{ErrorSemant} при этом сохраняется.


  


  Для данной серии переводов проставлена рубрика, характеризующая 


снижение качества НМП, так как  с~течением времени в~изначально 


допустимом варианте перевода появилась ошибка.


  


  Отметим, что при выборе рубрик <<Повышение\!/\!сни\-же\-ние качества 


НМП>> оценивается общая динамика качества, при этом наличие в~серии 


переводов допустимого варианта перевода (как в~табл.~\ref{t5-z}) не 


обязательно: в~табл.~\ref{t4-z} перевод остается ошибочным, но его качество 


повышается за счет уменьшения числа ошибок.


  


\section{Заключение}


  


  В статье приведено описание третьего этапа эксперимента по многократному 


фиксированию машинных переводов, полученных с~помощью системы НМП 


Google. Непосредственной целью данного этапа был количественный анализ 


результатов, полученных в~ходе исследования нестабильности НМП. В~итоге 


было выведено и~проанализировано соотношение категорий, характеризующих 


нестабильность НМП, и~видов ошибок НМП. Также описана динамика ошибок 


на протяжении всего эксперимента и~представлены данные распределения 


серий переводов по видам нестабильности.


  


  Полученные данные о~динамике ошибок свидетельствуют не только об 


актуальности изучения предельных возможностей обучения ИНС (что могут 


и~чего не могут делать компьютерные системы на основе ИНС~[17]), но 


и~о~необходимости индикаторной оценки результатов обучения при решении 


задач машинного первода. Разработанная система индикаторов и~их вычисленные значения 


уникальны. Предлагаемый подход к~индикаторной оценке результатов 


обучения ИНС может позиционироваться как новый аспект в~исследованиях 


качества машинного перевода (о~необходимости таких исследований см., например,~[18--20]).





{\small\frenchspacing


{%\baselineskip=10.8pt


\addcontentsline{toc}{section}{Литература}


\begin{thebibliography}{99}





\bibitem{1-z}


\Au{Fenstermacher K.\,D.} The tyranny of tacit knowledge: What artificial 


intelligence tells us about knowledge representation~/\!\!/ 38th Annual Hawaii 


Conference (International) on System Sciences Proceedings.~--- Washington, D.C., USA: 


IEEE Computer Society, 2005. Vol.~8. P.~243a. doi: 10.1109\!/\!HICSS.2005.620.


\bibitem{2-z}


\Au{Sako M.} Artificial intelligence and the future of professional work~/\!\!/ 


Commun. ACM, 2020. Vol.~63. No.\,4. P.~25--27.


\bibitem{3-z}


\Au{Denning P.\,J., Denning D.\,E.} Dilemmas of artificial intelligence~/\!\!/ 


Commun. ACM, 2020. Vol.~63. No.\,3. P.~22--24.


\bibitem{4-z}


\Au{Инькова-Манзотти О.\,Ю.} Коннекторы противопоставления во 


французском и~русском языках. Сопоставительное исследование.~--- М.: 


Информэлектро, 2001. 429~с.


\bibitem{9-z} %5


\Au{Зализняк А.\,А., Зацман И.\,М., Инькова О.\,Ю., Кружков~М.\,Г.} 


Надкорпусные базы данных как лингвистический ресурс~/\!\!/ Корпусная 


лингвистика: Тр. VII Междунар. конф.~--- СПб.: СПбГУ, 2015.  


С.~211--218.


\bibitem{8-z} %6


\Au{Дурново А.\,А., Зацман И.\,М., Лощилова Е.\,Ю.} Кросслингвистическая 


база данных для аннотирования логико-семантических отношений 


в~тексте~/\!\!/ Системы и~средства информатики, 2016. Т.~26. №\,4. С.~124--137.


\bibitem{6-z} %7


\Au{Зализняк А.\,А., Зацман И.\,М., Инькова О.\,Ю.} Надкорпусная база 


данных коннекторов: построение системы терминов~/\!\!/ Информатика и~её 


применения, 2017. Т.~11. Вып.~1. С.~100--108.


\bibitem{7-z} %8


\Au{Зацман И.\,М., Кружков М.\,Г.} Надкорпусная база данных коннекторов: 


развитие системы терминов проектирования~/\!\!/ Системы и~средства 


информатики, 2018. Т.~28. №\,4. С.~156--167.








\bibitem{10-z} %9


\Au{Бунтман Н.\,В., Гончаров А.\,А., Зацман И.\,М., Нуриев~В.\,А.} 


Количественный анализ результатов машинного перевода с~использованием 


надкорпусных баз данных~/\!\!/ Информатика и~её применения, 2018. Т.~12. 


Вып.~4.   С.~96--105.


\bibitem{5-z} %10


\Au{Егорова А.\,Ю., Зацман И.\,М., Мамонова О.\,С.} Надкорпусные базы 


данных в~лингвистических проектах~/\!\!/ Системы и~средства 


информатики, 2019. Т.~29. №\,3. С.~77--91.


\bibitem{11-z} %11


Национальный корпус русского языка. {\sf http://www.ruscorpora.ru}.


\bibitem{12-z}


\Au{Егорова А.\,Ю., Зацман И.\,М., Косарик В.\,В., Нуриев~В.\,А.} 


Нестабильность нейронного машинного перевода~/\!\!/ Системы и~средства 


информатики, 2020. Т.~30. №\,2. С.~124--135.


\bibitem{13-z}


\Au{Егорова А.\,Ю., Зацман И.\,М., Кружков М.\,Г., Нуриев~В.\,А.} Методика 


темпоральной оценки нестабильности машинного перевода~/\!\!/ Системы 


и~средства информатики, 2020. Т.~30. №\,3. С.~67--80.


\bibitem{14-z}


\Au{Гончаров А.\,А., Бунтман Н.\,В., Нуриев В.\,А.} Ошибки в~машинном 


переводе: проблемы классификации~/\!\!/ Системы и~средства информатики, 


2019. Т.~29. №\,3. С.~92--103.


\bibitem{15-z}


\Au{Инькова О.\,Ю.} Надкорпусная база данных как инструмент формальной 


вариативности коннекторов~/\!\!/ Компьютерная лингвистика 


и~интеллектуальные технологии: по мат-лам ежегодной Международ. конф.  


<<Диалог>>.~--- М.: РГГУ, 2018. Вып.~17(24). С.~240--253.


\bibitem{16-z}


\Au{Добровольский Д.\,О., Кретов А.\,А., Шаров С.\,А.} Корпус параллельных 


текстов: архитектура и~возможности использования~/\!\!/ Национальный 


корпус русского языка: 2003--2005~/ Отв. ред. В.\,А.~Плунгян.~--- М.: 


Индрик, 2005. С.~263--296.


\bibitem{17-z}


\Au{Brynjolfsson E., Mitchell T.} What can machine learning do? Workforce 


implications~/\!\!/ Science, 2017. Vol.~358. Iss.~6370. P.~1530--1534.





\bibitem{19-z} %18


\Au{Bowker L., Buitrago Ciro J.} Machine translation and global research: 


Towards improved machine translation literacy in the scholarly community.~--- 


Bingley, U.K.: Emerald Publishing, 2019. 111~p.


\bibitem{18-z} %19


\Au{Burchardt A., Lommel A., Macketanz V.} A~new deal for translation 


quality~/\!\!/ Universal Access Inf., 2020. doi: 10.1007\!/\!s10209-020-00736-5.


\bibitem{20-z}


\Au{Koehn Ph.} Neural machine translation.~--- New York, NY, USA: Cambridge 


University Press, 2020. 408~p.





\end{thebibliography}


} }








\vspace*{-3pt}





\hfill{\small\textit{Поступила в~редакцию 14.09.20}}














\vspace*{6pt}





\hrule





\vspace*{2pt}





%\newpage





%\vspace*{-28pt}





\hrule





\vspace*{8pt}





\def\tit{MACHINE TRANSLATION: INDICATOR-BASED EVALUATION 


OF~TRAINING PROGRESS IN~NEURAL PROCESSING}





\def\titkol{Machine translation: Indicator-based evaluation 


of training progress in neural processing}











\def\aut{A.\,Yu.~Egorova, I.\,M.~Zatsman, M.\,G.~Kruzhkov, and 


V.\,A.~Nuriev }





\def\autkol{A.\,Yu.~Egorova, I.\,M.~Zatsman, M.\,G.~Kruzhkov, and 


V.\,A.~Nuriev}





\titel{\tit}{\aut}{\autkol}{\titkol}





\vspace*{-11pt}











\noindent


Institute of Informatics Problems, Federal Research Center ``Computer Science and 


Control'' of the Russian Academy of Sciences,  


44-2~Vavilov Str., Moscow 119333, Russian Federation





\def\leftfootline{\small{\textbf{\thepage}


\hfill Sistemy i Sredstva Informatiki~---


Systems and Means of Informatics 2020 vol~30 no\,4}


}%


 \def\rightfootline{\small{Sistemy i Sredstva 


Informatiki~---


 Systems and Means of Informatics 2020 vol~30 no\,4


\hfill \textbf{\thepage}}}


  


  


\vspace*{3pt}








\Abste{The paper presents data collected while observing training progress of a 


neural machine translation (NMT) engine. The observed training progress received 


qualitative evaluation based on a set of indicators. Two hundred and fifty text fragments in Russian 


were used as experimental material for the study. For the duration of one 


year, every month these fragments were translated into French using the publicly 


available Google's NMT engine. The produced translations were recorded and 


annotated by language experts in a supracorpora database which resulted in a series 


of 12 annotated translations for each of the 250 Russian fragments. The annotations 


include labels of translation errors which enables researchers to determine the NMT 


instability types according to the changes of translation quality or lack thereof. The goal 


of this paper is to describe the newly developed indicator-based approach and to 


provide an example of its application to evaluation of a neural network training 


progress.}





\KWE{neural machine translation; instability of machine translation; indicator-based 


evaluation; linguistic annotation; instability types}





\DOI{10.14357\!/\!08696527200412} 








%\vspace*{-24pt}





%\Ack


%\noindent








\vspace*{-12pt}





\renewcommand{\bibname}{\protect\rmfamily References}


%\renewcommand{\bibname}{\large\protect\rm References}





{\small\frenchspacing


{%\baselineskip=10.8pt


\addcontentsline{toc}{section}{References}


\begin{thebibliography}{99}





\bibitem{1-z-1}


\Aue{Fenstermacher, K.\,D.} 2005. The tyranny of tacit knowledge: What artificial 


intelligence tells us about knowledge representation. \textit{38th Annual Hawaii 


Conference (International) on System Sciences Proceedings}. Washington, D.C.:


IEEE Computer Society. 8:243a.


doi: 10.1109\!/\!HICSS.2005.620.


\bibitem{2-z-1}


\Aue{Sako, M.} 2020. Artificial intelligence and the future of professional work. 


\textit{Commun. ACM}   63(4):25--27.


\bibitem{3-z-1}


\Aue{Denning, P.\,J., and D.\,E. Denning.} 2020. Dilemmas of artificial intelligence. 


\textit{Commun. ACM} 63(3):22--24.


\bibitem{4-z-1}


\Aue{Inkova-Manzotti, O.\,Yu.} 2001. \textit{Konnektory protivopostavleniya vo 


frantsuzskom i~russkom yazykakh. Sopostavitel'noe issledovanie} [Connectors of 


opposition in French and Russian: A~comparative study]. Moscow: Informelektro. 


429~p.


\bibitem{9-z-1} %5


\Aue{Zaliznyak, A.\,A., I.\,M. Zatsman, O.\,Yu. Inkova, and M.\,G.~Kruzhkov.} 


2015. Nadkorpusnye bazy dannykh kak lingvisticheskiy resurs [Supracorpora 


databases as linguistic resource]. \textit{Tr. VII Mezhdunar. Konf. ``Korpusnaya 


lingvistika''} [7th Conference (International) ``Corpus Linguistics'' Proceedings]. 


St.\ Petersburg: SPbSU. 211--218.





\bibitem{8-z-1} %6


\Aue{Durnovo, A.\,A., I.\,M. Zatsman, and E.\,Yu.~Loshchilova.} 2016. 


Krosslingvi\-sti\-che\-skaya baza dannykh dlya annotirovaniya logiko-semanticheskikh 


otnosheniy v~tekste [Cross-lingual database for annotating logical-semantic relations 


in the text]. \textit{Sistemy i~Sredstva Informatiki~--- Systems and Means of 


Informatics} 26(4):124--137. 





\bibitem{6-z-1} %7


\Aue{Zaliznyak, A.\,A., I.\,M. Zatsman, and O.\,Yu. Inkova.} 2017. Nadkorpusnaya 


baza dannykh konnektorov: postroenie sistemy terminov [Supracorpora database on 


connectives: Term system development]. \textit{Informatika i~ee  


Primeneniya~--- Inform. Appl.} 11(1):100--108.


\bibitem{7-z-1} %8


\Aue{Zatsman, I.\,M., and M.\,G. Kruzhkov.} 2018. Nadkorpusnaya baza dannykh 


konnektorov: razvitie sistemy terminov proektirovaniya [Supracorpora database of 


connectives: Design-oriented evolution of the term system]. \textit{Sistemy 


i~Sredstva Informatiki~--- Systems and Means of Informatics} 28(4):156--167.








\bibitem{10-z-1} %9


\Aue{Buntman, N.\,V., A.\,A. Goncharov, I.\,M. Zatsman, and V.\, A.~Nuriev.} 


2018. Ko\-li\-chest\-ven\-nyy analiz rezul'tatov mashinnogo perevoda s~ispol'zovaniem 


nadkorpusnykh baz dannykh [Using supracorpora databases for quantitative analysis 


of machine translations]. \textit{Informatika i ee Primeneniya~--- Inform. Appl.} 


12(4):96--105. 


\bibitem{5-z-1} %10


\Aue{Egorova, A.\,Yu., I.\,M. Zatsman, and O.\,S. Mamonova.} 2019. Nadkorpusnye 


bazy dannykh v~lingvisticheskikh proektakh [Supracorpora databases in linguistic 


projects]. \textit{Sistemy i~Sredstva Informatiki~--- Systems and Means of 


Informatics} 29(3):77--91.


\bibitem{11-z-1}


Natsional'nyy korpus russkogo yazyka [Russian National Corpus]. Available at: {\sf 


http://www.ruscorpora.ru/} (accessed September~11, 2020).


\bibitem{12-z-1}


\Aue{Egorova, A.\,Yu., I.\,M. Zatsman, V.\,V. Kosarik, and V.\,A.~Nuriev.} 2020. 


Nestabil'nost' neyronnogo mashinnogo perevoda [Instability of neural machine 


translation]. \textit{Sistemy i~Sredstva Informatiki~--- Systems and Means of 


Informatics} 30(2):124--135.


\bibitem{13-z-1}


\Aue{Egorova, A.\,Yu., I.\,M. Zatsman, M.\,G. Kruzhkov, and V.\,A.~Nuriev.} 2020. 


Metodika temporal'noy otsenki nestabil'nosti mashinnogo perevoda [The technique 


allowing for temporal estimation of the machine translation instability]. 


\textit{Sistemy i~Sredstva Informatiki~--- Systems and Means of Informatics} 


30(3):67--80.


\bibitem{14-z-1}


\Aue{Goncharov, A.\,A., N.\,V. Buntman, and V.\,A. Nuriev.} 2019. Oshibki 


v~mashinnom perevode: problemy klassifikatsii [Machine translation errors: 


Problems of classification]. \textit{Sistemy i~Sredstva Informatiki~--- Systems and 


Means of Informatics} 29(3):92--103.


\bibitem{15-z-1}


\Aue{Inkova, O.\,Yu.} 2018. Nadkorpusnaya baza dannykh kak instrument 


formal'noy variativnosti konnektorov [Supracorpora database as an instrument of the 


study of the formal variability of connectives]. \textit{Komp'yuternaya lingvistika 


i~intellektual'nye tekhnologii: po mat-lam ezhegodnoy Mezhdunar. konf. ``Dialog''} 


[Computer Linguistic and Intellectual Technologies: Conference (International) 


``Dialog'' Proceedings]. Moscow. 17(24):240--253.


\bibitem{16-z-1}


\Aue{Dobrovol'skiy, D.\,O., A.\,A. Kretov, and S.\,A. Sharov.} 2005. Korpus 


parallel'nykh tekstov: arkhitektura i~vozmozhnosti ispol'zovaniya [Corpus of 


parallel texts: Architecture and applications]. \textit{Natsional'nyy korpus russkogo 


yazyka: 2003--2005} [Russian National Corpus: 2003--2005]. Moscow: Indrik. 263--296.


\bibitem{17-z-1}


\Aue{Brynjolfsson, E., and T. Mitchell.} 2017. What can machine learning do? 


Workforce implications. \textit{Science} 358(6370):1530--1534.





\bibitem{19-z-1} %18


\Aue{Bowker, L., and J.\,B. Ciro.} 2019. \textit{Machine translation and global 


research: Towards improved machine translation literacy in the scholarly 


community}. Bingley, U.K.: Emerald Publishing. 111~p.


\bibitem{18-z-1} %19


\Aue{Burchardt, A., A. Lommel, and V.~Macketanz.} 2020. A new deal for 


translation quality. \textit{Universal Access Inf.}


doi: 10.1007/s10209-020-00736-5.





\bibitem{20-z-1}


\Aue{Koehn, Ph.} 2020. \textit{Neural machine translation}. New York, NY: 


Cambridge University Press. 408~p.


\end{thebibliography}


} }





%\vspace*{-3pt}





\hfill{\small\textit{Received September 14, 2020}} 





%\vspace*{-16pt}














\Contr





\noindent


\textbf{Egorova Anna Yu.} (b.\ 1991)~--- junior scientist, Institute of Informatics 


Problems, Federal Research Center ``Computer Science and Control'' of the Russian 


Academy of Sciences, 44-2~Vavilov Str., Moscow 119333, Russian Federation; 


\mbox{ann.shurova@gmail.com}





\vspace*{3pt}





\noindent


\textbf{Zatsman Igor M.} (b.\ 1952)~--- Doctor of Science in technology, Head of 


Department, Institute of Informatics Problems, Federal Research Center ``Computer 


Science and Control'' of the Russian Academy of Sciences, 44-2~Vavilov Str., 


Moscow 119333, Russian Federation; \mbox{izatsman@yandex.ru}





\vspace*{3pt}


\noindent


\textbf{Kruzhkov Mikhail G.} (b.\ 1975)~--- senior scientist, Institute of Informatics 


Problems, Federal Research Center ``Computer Science and Control'' of the Russian 


Academy of Sciences, 44-2~Vavilov Str., Moscow 119333, Russian Federation; 


\mbox{magnit75@yandex.ru}





\vspace*{3pt}





\noindent


\textbf{Nuriev Vitaly A.} (b.\ 1980)~--- Candidate of Science (PhD) in philology, 


leading scientist, Institute of Informatics Problems, Federal Research Center 


``Computer Science and Control'' of the Russian Academy of Sciences, 


44-2~Vavilov Str., Moscow 119333, Russian Federation; \mbox{nurieff.v@gmail.com}





\label{end\stat}





\renewcommand{\bibname}{\protect\rm Литература} 


